Sobre los 496 pares de ciudades, 385 tienen dos caminos simples y 111 tienen uno solo. Del total, 374 pares tienen paso obligatorio (un nodo intermedio presente en todos sus caminos), 105 tienen atajo (dos caminos sin ningún nodo intermedio en común) y 17 son de camino único, sin nodo intermedio. La red tiene 10 puntos de articulación, nodos cuya eliminación la desconecta: Bogotá, Fusagasugá, Girardot, Honda, Puerto López, Sáchica, Tunja, Villavicencio, Villeta y Ye de Granada.
